\subsection{RELATIONS BETWEEN RESTRICTION AND EXTENSION OF THE RING OF SCALARS}

Let \(\rho: \mathbf{A} \to \mathbf{B}\) be a ring homomorphism. For every left A-module E, a canonical A-linear mapping

\[
\phi_ {\mathbf {E}}: \mathbf {E} \rightarrow \rho_ {*} (\rho^ {*} (\mathbf {E}))\tag{10}
\]

was defined in no. 1 such that \(\phi_{\mathbb{E}}(x) = 1 \otimes x\) . We consider now a left B-module F and apply Proposition 1 (no. 1) to the A homomorphism \(\mathbf{l}_{\rho_{*}(\mathbf{F})}: \rho_{*}(\mathbf{F}) \to \rho_{*}(\mathbf{F})\) : we obtain a B-linear mapping

\[
\psi_ {\mathbf {F}} \colon \rho^ {*} (\rho_ {*} (\mathbf {F})) \rightarrow \mathbf {F}\tag{11}
\]

equal to \(\delta(1_{\rho_*(\mathbb{F})})\) and such therefore that, for all \(y \in \mathbb{F}\) and all \(\beta \in \mathbb{B}\) , \(\psi_{\mathbb{F}}(\beta \otimes y) = \beta y\) .

PROPOSITION 5. Let E be a left A-module and F a left B-module; the composite mappings

\[
\rho^ {*} (\mathrm{E}) \xrightarrow [ \rho^ {*} (\phi_ {\mathrm{E}}) ]{} \rho^ {*} (\rho_ {*} (\rho^ {*} (\mathrm{E}))) \xrightarrow [ \psi_ {\rho^ {*} (\mathrm{E})} ]{} \rho^ {*} (\mathrm{E})\tag{12}
\]

\[
\rho_ {*} (\mathrm{F}) \xrightarrow [ \phi_ {\rho_ {*} (\mathrm{F})} ]{} \rho_ {*} (\rho^ {*} (\rho_ {*} (\mathrm{F}))) \xrightarrow [ \rho_ {*} (\psi_ {\mathrm{F}}) ]{} \rho_ {*} (\mathrm{F})\tag{13}
\]

are respectively equal to the identity mappings of \(\rho^{*}(\mathbf{E})\) and \(\rho_{*}(\mathbf{F})\) .

We give the proof, for example, for (12); for all \(x \in \mathbf{E}\) , the mapping \(\rho^{*}(\phi_{\mathbf{E}})\) maps \(1 \otimes x\) to the element \(1 \otimes (1 \otimes x)\) and the mapping \(\psi_{\rho^{*}(\mathbf{E})}\) maps \(1 \otimes (1 \otimes x)\) to the element \(1 \otimes x\) ; the conclusion follows from the fact that the elements of the form \(1 \otimes x\) generate the B-module \(\rho^{*}(\mathbf{E})\) ; the proof is even simpler for (13).

COROLLARY. The mappings \(\rho^{*}(\phi_{\mathbf{E}})\) and \(\phi_{\rho^{*}(\mathbf{F})}\) are injective and respectively identify \(\rho^{*}(\mathbf{E})\) with a direct factor of \(\rho^{*}(\rho_{*}(\rho^{*}(\mathbf{E})))\) and \(\rho_{*}(\mathbf{F})\) with a direct factor of \(\rho_{*}(\rho^{*}(\rho_{*}(\mathbf{F})))\) .

This is a consequence of Proposition 5 and §1, no. 9, Corollary 2 to Proposition 15.

PROPOSITION 6. Let E be a left A-module and F a right B-module. There exists one and only one Z-homomorphism

\[
\rho_ {*} (\mathrm{F}) \otimes_ {\mathrm{A}} \mathrm{E} \rightarrow \mathrm{F} \otimes_ {\mathrm{B}} \rho^ {*} (\mathrm{E})\tag{14}
\]

mapping \(y \otimes x\) to \(y \otimes (1 \otimes x)\) for all \(x \in E\) and all \(y \in F\) and this homomorphism is bijective.

By definition the right hand side of (14) is \(F \otimes_{B} (B \otimes_{A} E)\) , where B is considered as a (B, A)-bimodule, and there is a canonical Z-isomorphism \((\mathbf{F} \otimes_{\mathbf{B}} \mathbf{B}) \otimes_{\mathbf{A}} \mathbf{E} \to \mathbf{F} \otimes_{\mathbf{B}} (\mathbf{B} \otimes_{\mathbf{A}} \mathbf{E})\) defined in §3, no. 8, Proposition 8; on the other hand, the canonical isomorphism \(F \to F \otimes_{B} B\) of §3, no. 4, Proposition 4 is an isomorphism for the right A-module structures on the two sides, defined by \(\rho\) . Whence the desired isomorphism.

When A and B are commutative, the isomorphism (14) is an A-module isomorphism

\[
\rho_ {*} (\mathrm{F}) \otimes_ {\mathrm{A}} \mathrm{E} \rightarrow \rho_ {*} (\mathrm{F} \otimes_ {\mathrm{B}} \rho^ {*} (\mathrm{E})).
\]

